\chapter{LITERATURE REVIEW}
\section{Overview}
Intersections are key junctions in the road networks where 
multiple roads meet or cross. Their design and traffic control 
methods play a vital role in ensuring smooth traffic flow, safety, 
and overall efficiency. Traffic management at intersections has 
been a complex issue in transportation systems for a long time. 
Several real-time signal control strategies have been implemented, 
but have not been able to effectively solve this problem completely 
\cite{tomar2022}.

Broadly, intersections are categorized into signalized and 
unsignalized types based on the control mechanism used to manage 
traffic movement.

\section{Signalized Intersection}
A signalized intersection uses traffic signals to control the 
movement of vehicles and pedestrians. These signals show green, 
yellow, and red lights to indicate when to proceed, slow down, or 
stop. The signals operate on a timed cycle, assigning specific 
green light durations to different directions to maintain smooth 
and orderly traffic flow. \cite{duraku2024}. Traffic 
signals help organize the flow of vehicles and people at intersections, 
making sure everyone gets their turn to move safely. By giving clear 
directions on when to stop and go, they help prevent traffic jams 
and reduce the chances of accidents. This makes the roads safer 
and more efficient for both drivers and pedestrians. “The traffic 
signal timing plan is the key that assigns the required time based 
on the travel demand at the intersection and keeps the cycle length 
to a minimum green time”. The signal controller works as the time 
setting device for the specific intersection \cite{fhwa2008}. Traffic 
flow movements are controlled by traffic signals without any 
conflict points in the signalized intersections. The movements 
are indicated by green, yellow, and red signals for the vehicle 
and pedestrian. The road users in the traffic include passenger 
cars, heavy goods vehicles, large goods vehicles, cycles, pedestrians, 
public transportation, and motorbikes \cite{eom2020}. Traffic s
ignals are electrically operated control devices to avoid conflicts 
between vehicles and pedestrians \cite{ryder2005}. The 
three main concepts that describe the traffic signal are the cycle, 
cycle length, and phase. The cycle is the complete sequence of all 
the phases. Cycle length is the duration of one complete cycle. 
The phase is the time interval during which a specific group of 
vehicle traffic movements is permitted or restricted at the signalized 
intersection. The signal control cycle has different phases, with 
each phase representing a combination of permitted movements for 
multiple lanes receiving the right of way simultaneously \cite{hcm2000}. 
The basic control is applied through FTSC, and ATSC strategies include 
the control of FATSC (Full Actuated Traffic Signal Control), SATSC 
(Semi Actuated Traffic Signal Control), and Adaptive Traffic Signal 
Control (ADTSC). These signal systems adjust according to traffic needs 
using data from detectors \cite{duraku2024}

\section{Pretimed Signal Control}

Pretimed control, also known as fixed time signal control, is a 
signal control system in which the cycle length, phase plan, and 
phase times are preset to repeat continuously. In this system, a 
preset time is given to each movement every cycle, regardless of 
changes in traffic conditions. Pretimed signal control operates on 
a fixed schedule, where each movement (e.g., vehicular or pedestrian 
phases) is allocated a predetermined amount of time within a 
signal cycle regardless of real-time traffic demand \cite{fhwa2007}. 
Based on historic data at the intersection, the signal timing values 
were calculated, and these values are programmed into the pre-timed 
controllers. The pre-timed signal control has many advantages. 
For example, since both the start and end of green are predictable, 
it can be used for efficient coordination with adjacent pre-timed 
signals. Also, as this system does not require detectors, its 
operations are immune to problems that arise due to the failure 
of detectors. But on the other hand, fixed signal time control 
becomes inefficient at isolated intersections where traffic 
arrivals are very low. Pre-timed control is ideally suited to 
closely spaced intersections, mostly found in downtown areas. 
Also, this system is a great option for those intersections where 
a maximum of three phases is required \cite{urbanik2015}. In 
the absence of traffic fluctuations in traffic demands, optimum 
fixed timed management (OFTM) is widely used in isolated intersection 
approaches to manage traffic flow. Fluctuations in traffic demand 
decrease the efficiency of OFTM, which causes longer queue lengths, 
delay, fuel consumption, and exhaust emissions  \cite{grether2011}. 
Changing the timing of traffic signals is a cheap and simple way 
to improve traffic flow and reduce congestion. \cite{fhwa2007}. But 
traditional fixed-time traffic signals operate on preset schedules 
and often fail to respond effectively to changing traffic volumes, 
especially during peak hours \cite{costa2018}.

\section{Actuated Signal Control}

Whenever there are short-term or long-term fluctuations in traffic 
demand at some intersection approaches, the effectiveness and 
efficiency of the Optimum Fixed Time Management (OFTM) system 
might decrease or can be adversely affected\cite{grether2011}. 
Traffic management systems based on any created algorithm are 
referred to as vehicle-actuated signal management systems. A 
vehicle-actuated management system determines the green and red 
signal timings according to the current dynamics of vehicle arrivals 
and queuing information obtained by detectors placed on the lanes 
in the intersection approaches\cite{akcelik1994}. One of the applications 
of Intelligent Transportation Systems (ITS) is vehicle-actuated (traffic- actuated) 
management systems. The effectiveness of these systems depends 
heavily on the control logic. The selection of the most appropriate 
control parameters for the designed control logic determines the 
effectiveness and functions of the actuated systems \cite{cakici2021}. 
Vehicle-actuated controllers are also a type of adaptive system, but 
the term “adaptive” usually refers to an area traffic control 
system (ATSC) where signal timings at various junctions are 
changed together based on that real time traffic conditions
\cite{ravikumar2012}. An ATCS is a complex system that involves 
extensive surveillance, such as pavement loop detectors, and 
infrastructure to connect with central and local controllers. 
However, these systems are highly efficient, as they include 
algorithms that adjust a signal’s split, offset, phase length, 
and phase sequences to control and minimize delays, ultimately 
reducing the number of stops \cite{stevanovic2010}. The analysis 
results show that the developed vehicle-actuated management system 
can effectively adapt to fluctuations in traffic demand at signalized 
intersection approaches.

The system significantly reduces average vehicle delays, fuel 
consumption, and exhaust emissions. Implementation of a vehicle-actuated 
management system can greatly improve the performance of an intersection 
in the presence of fluctuating or varying traffic demand
\cite{cakici2021}. The actuated signal controller has detectors 
on all the approach lanes. When a detector detects the vehicle, 
the information is transferred to the controller, and the controller 
will give the service needed for the road users. The current signal 
tage will reach its minimum green time if they do not have more 
vehicles, and then it will change to the next stage to give the 
right of way to the vehicles. Detectors are placed on the lanes 
to communicate with the controller, and they should be properly 
numbered. The detector numbering depends on the phase numbering 
or the type of signal controller used. Fully actuated signal systems 
adjust the traffic lights based on real-time data collected from 
all lanes at an intersection, making these systems the most flexible. 
Semi-actuated systems only collect data from some lanes, so they 
only gather partial demand information and make decisions based 
on that data. Fully actuated control works best at isolated intersections 
where there are fluctuations in traffic demand throughout the day 
\cite{urbanik2015}. The maximum green time should be set for 
the actuated signals to determine the green signal splits for the 
intersection. If the maximum green time is set too high, then the 
high-demand lanes would start to affect the operations on the other 
approach lanes. The minimum green time is the amount of green time 
required for the vehicle to clear the signalized intersections; it 
should be enough for the vehicle to pass the intersection, otherwise 
it will affect the efficiency of the intersection (Sinowatcher, n.d.).

In a semi-actuated control system, the design logic control dynamically 
adjusts signal timings by determining the green light duration for 
the secondary road based on the arriving traffic flow. When the detector 
detects no vehicles or the absence of traffic, the control logic eliminates 
the green interval and moves to the next phase, which greatly reduces 
cycle time and delay. This dynamic adjustment of the switching signal operates 
instantly using an 'if-then' conditional logic to optimize traffic signal 
performance  \cite{duraku2024}. Vehicle-actuated signal systems were 
developed to improve upon fixed-time systems by detecting vehicles approaching 
an intersection and adjusting green time accordingly.

Although more efficient than fixed-time systems, vehicle-actuated 
signals are still limited in their ability to coordinate across 
multiple intersections (Brown \& Patel, 2019).

The cost for both the initial set-up and the maintenance is the 
major disadvantage of fully actuated control. The high cost is 
due to the amount of detection required. Also, there are chances 
for a higher percentage of vehicles to stop since the green time 
is not held for upstream traffic. Another disadvantage is that, 
if the traffic patterns are regular, the extra benefit obtained 
from the addition of local actuation is minimal or can be zero 
\cite{fhwa2008}. Using a fuzzy logic-based vehicle actuated management 
system reduced delays by 10% (Trabia et al). The green time extension 
investigation for vehicle-actuated signal control \cite{wang2018}. 
The green and red signal timings at each cycle were determined using 
global optimization and complementarity for vehicle-actuated management 
systems (Ribeiro and Simoes et al). At the end of the analysis, it 
was concluded that a fully-actuated management system could provide 
better and more encouraging results than the fixed- time and semi-actuated 
management systems.

\section{Adaptive Signal Control}

Adaptive Traffic Control Systems (ATCSs) change signal timings in 
real time based on how much traffic is on the road and how the 
system is working. These systems use smart rules to change things 
like how long lights stay green, when they start, and the order 
of signals to reduce delays and stops. They need a lot of monitoring, 
often using sensors under the road, and special equipment to connect 
with control centers\cite{stevanovic2010}. Adaptive traffic control 
systems are developed to solve the limitations of pre-timed control 
by instantly adapting to field conditions, such as changing traffic 
signal timings based on the traffic present and fluctuations in 
demand or movement\cite{martin2008}. Various adaptive traffic control 
systems like SCOOT, SCATS, and OPAC optimize signal timings based 
on real-time and predicted traffic demand to improve flow through 
coordinated green waves. However, the installation and maintenance 
costs of these systems are very high, and they also require complex 
infrastructure. As a result, they are not widely implemented and 
are built in only a few cities around the world \cite{tomar2022}.  